% Options for packages loaded elsewhere
% Options for packages loaded elsewhere
\PassOptionsToPackage{unicode}{hyperref}
\PassOptionsToPackage{hyphens}{url}
\PassOptionsToPackage{dvipsnames,svgnames,x11names}{xcolor}
%
\documentclass[
  11pt,
]{article}
\usepackage{xcolor}
\usepackage[margin=1in]{geometry}
\usepackage{amsmath,amssymb}
\setcounter{secnumdepth}{-\maxdimen} % remove section numbering
\usepackage{iftex}
\ifPDFTeX
  \usepackage[T1]{fontenc}
  \usepackage[utf8]{inputenc}
  \usepackage{textcomp} % provide euro and other symbols
\else % if luatex or xetex
  \usepackage{unicode-math} % this also loads fontspec
  \defaultfontfeatures{Scale=MatchLowercase}
  \defaultfontfeatures[\rmfamily]{Ligatures=TeX,Scale=1}
\fi
\usepackage{lmodern}
\ifPDFTeX\else
  % xetex/luatex font selection
\fi
% Use upquote if available, for straight quotes in verbatim environments
\IfFileExists{upquote.sty}{\usepackage{upquote}}{}
\IfFileExists{microtype.sty}{% use microtype if available
  \usepackage[]{microtype}
  \UseMicrotypeSet[protrusion]{basicmath} % disable protrusion for tt fonts
}{}
\makeatletter
\@ifundefined{KOMAClassName}{% if non-KOMA class
  \IfFileExists{parskip.sty}{%
    \usepackage{parskip}
  }{% else
    \setlength{\parindent}{0pt}
    \setlength{\parskip}{6pt plus 2pt minus 1pt}}
}{% if KOMA class
  \KOMAoptions{parskip=half}}
\makeatother
% Make \paragraph and \subparagraph free-standing
\makeatletter
\ifx\paragraph\undefined\else
  \let\oldparagraph\paragraph
  \renewcommand{\paragraph}{
    \@ifstar
      \xxxParagraphStar
      \xxxParagraphNoStar
  }
  \newcommand{\xxxParagraphStar}[1]{\oldparagraph*{#1}\mbox{}}
  \newcommand{\xxxParagraphNoStar}[1]{\oldparagraph{#1}\mbox{}}
\fi
\ifx\subparagraph\undefined\else
  \let\oldsubparagraph\subparagraph
  \renewcommand{\subparagraph}{
    \@ifstar
      \xxxSubParagraphStar
      \xxxSubParagraphNoStar
  }
  \newcommand{\xxxSubParagraphStar}[1]{\oldsubparagraph*{#1}\mbox{}}
  \newcommand{\xxxSubParagraphNoStar}[1]{\oldsubparagraph{#1}\mbox{}}
\fi
\makeatother


\usepackage{longtable,booktabs,array}
\usepackage{calc} % for calculating minipage widths
% Correct order of tables after \paragraph or \subparagraph
\usepackage{etoolbox}
\makeatletter
\patchcmd\longtable{\par}{\if@noskipsec\mbox{}\fi\par}{}{}
\makeatother
% Allow footnotes in longtable head/foot
\IfFileExists{footnotehyper.sty}{\usepackage{footnotehyper}}{\usepackage{footnote}}
\makesavenoteenv{longtable}
\usepackage{graphicx}
\makeatletter
\newsavebox\pandoc@box
\newcommand*\pandocbounded[1]{% scales image to fit in text height/width
  \sbox\pandoc@box{#1}%
  \Gscale@div\@tempa{\textheight}{\dimexpr\ht\pandoc@box+\dp\pandoc@box\relax}%
  \Gscale@div\@tempb{\linewidth}{\wd\pandoc@box}%
  \ifdim\@tempb\p@<\@tempa\p@\let\@tempa\@tempb\fi% select the smaller of both
  \ifdim\@tempa\p@<\p@\scalebox{\@tempa}{\usebox\pandoc@box}%
  \else\usebox{\pandoc@box}%
  \fi%
}
% Set default figure placement to htbp
\def\fps@figure{htbp}
\makeatother





\setlength{\emergencystretch}{3em} % prevent overfull lines

\providecommand{\tightlist}{%
  \setlength{\itemsep}{0pt}\setlength{\parskip}{0pt}}



 


% ============================================================
%  Econ 502 — Problem Set / Handout Theme
% ============================================================

% ---------- packages ----------
\usepackage{enumitem}
\usepackage{tikz}
\usepackage{fancyhdr}
\usepackage{titlesec}
\usepackage{tcolorbox}
\usepackage{booktabs}

\tcbuselibrary{skins, breakable}

% ---------- colors ----------
\definecolor{csufblue}{HTML}{1E5288}
\definecolor{csufnavy}{HTML}{00274C}
\definecolor{accentteal}{HTML}{16A085}
\definecolor{lightgray}{HTML}{F5F5F5}
\definecolor{medgray}{HTML}{E0E0E0}
\definecolor{darktext}{HTML}{1A1A1A}

% ---------- page style ----------
\pagestyle{fancy}
\fancyhf{}
\renewcommand{\headrulewidth}{0.8pt}
\renewcommand{\headrule}{\hbox to\headwidth{\color{csufblue}\leaders\hrule height \headrulewidth\hfill}}
\fancyhead[L]{\small\color{csufblue}\textbf{Econ 502: Advanced Microeconomics}}
\fancyhead[R]{\small\color{csufblue}\thepage}
\fancyfoot{}

% first page: colored footer only
\fancypagestyle{plain}{%
  \fancyhf{}
  \renewcommand{\headrulewidth}{0pt}
  \fancyfoot[C]{\small\color{csufblue}\thepage}
}

% ---------- title block ----------
% Hook into Pandoc's \maketitle: add ruled lines around the title
\makeatletter
\renewcommand{\maketitle}{%
  \thispagestyle{plain}
  \begin{center}
    {\color{csufblue}\rule{\textwidth}{1.5pt}}\par
    \vspace{10pt}
    {\LARGE\bfseries\color{csufnavy}\@title\par}
    \vspace{8pt}
    {\color{csufblue}\rule{\textwidth}{1.5pt}}
  \end{center}
  \vspace{12pt}
}
\makeatother

% ---------- section headings ----------
\titleformat{\section}
  {\Large\bfseries\color{csufnavy}}
  {}{0pt}{}
  [\vspace{2pt}{\color{csufblue}\titlerule[0.8pt]}]

\titleformat{\subsection}
  {\large\bfseries\color{csufblue}}
  {}{0pt}{}

\titleformat{\subsubsection}
  {\normalsize\bfseries\color{accentteal}}
  {}{0pt}{}

\titlespacing*{\section}{0pt}{16pt}{8pt}
\titlespacing*{\subsection}{0pt}{12pt}{4pt}
\titlespacing*{\subsubsection}{0pt}{8pt}{4pt}

% ---------- tcolorbox environments ----------

% Answer/solution box (blue)
\newtcolorbox{answer}{
  colback   = csufblue!5,
  colframe  = csufblue,
  arc       = 3pt,
  boxrule   = 0.8pt,
  left      = 8pt,
  right     = 8pt,
  top       = 6pt,
  bottom    = 6pt,
  breakable,
}

% Key result box (teal, with optional title)
\newtcolorbox{keyresult}[1][]{
  colback      = accentteal!5,
  colframe     = accentteal,
  arc          = 3pt,
  boxrule      = 0.8pt,
  left         = 8pt,
  right        = 8pt,
  top          = 6pt,
  bottom       = 6pt,
  fonttitle    = \bfseries\color{white},
  coltitle     = white,
  colbacktitle = accentteal,
  title        = {#1},
  breakable,
}

% Note box (light gray)
\newtcolorbox{note}{
  colback  = lightgray,
  colframe = medgray,
  arc      = 3pt,
  boxrule  = 0.6pt,
  left     = 8pt,
  right    = 8pt,
  top      = 6pt,
  bottom   = 6pt,
  breakable,
}

% ---------- list styling ----------
\setlist[enumerate, 1]{label=(\alph*), leftmargin=*, topsep=4pt, itemsep=4pt}
\setlist[enumerate, 2]{label=(\roman*), leftmargin=*, topsep=2pt, itemsep=2pt}
\setlist[itemize]{leftmargin=*, topsep=4pt, itemsep=3pt}

% ---------- boxed answers ----------
% Colored \boxed command
\renewcommand{\boxed}[1]{\fcolorbox{csufblue}{csufblue!4}{$\displaystyle #1$}}

% ---------- hyperlink colors ----------
\AtBeginDocument{%
  \hypersetup{
    colorlinks = true,
    linkcolor  = csufblue,
    urlcolor   = accentteal,
    citecolor  = csufnavy,
  }%
}
\makeatletter
\@ifpackageloaded{caption}{}{\usepackage{caption}}
\AtBeginDocument{%
\ifdefined\contentsname
  \renewcommand*\contentsname{Table of contents}
\else
  \newcommand\contentsname{Table of contents}
\fi
\ifdefined\listfigurename
  \renewcommand*\listfigurename{List of Figures}
\else
  \newcommand\listfigurename{List of Figures}
\fi
\ifdefined\listtablename
  \renewcommand*\listtablename{List of Tables}
\else
  \newcommand\listtablename{List of Tables}
\fi
\ifdefined\figurename
  \renewcommand*\figurename{Figure}
\else
  \newcommand\figurename{Figure}
\fi
\ifdefined\tablename
  \renewcommand*\tablename{Table}
\else
  \newcommand\tablename{Table}
\fi
}
\@ifpackageloaded{float}{}{\usepackage{float}}
\floatstyle{ruled}
\@ifundefined{c@chapter}{\newfloat{codelisting}{h}{lop}}{\newfloat{codelisting}{h}{lop}[chapter]}
\floatname{codelisting}{Listing}
\newcommand*\listoflistings{\listof{codelisting}{List of Listings}}
\makeatother
\makeatletter
\makeatother
\makeatletter
\@ifpackageloaded{caption}{}{\usepackage{caption}}
\@ifpackageloaded{subcaption}{}{\usepackage{subcaption}}
\makeatother
\usepackage{bookmark}
\IfFileExists{xurl.sty}{\usepackage{xurl}}{} % add URL line breaks if available
\urlstyle{same}
\hypersetup{
  pdftitle={CES Production Function},
  colorlinks=true,
  linkcolor={blue},
  filecolor={Maroon},
  citecolor={Blue},
  urlcolor={Blue},
  pdfcreator={LaTeX via pandoc}}


\title{CES Production Function}
\usepackage{etoolbox}
\makeatletter
\providecommand{\subtitle}[1]{% add subtitle to \maketitle
  \apptocmd{\@title}{\par {\large #1 \par}}{}{}
}
\makeatother
\subtitle{Econ 502: Advanced Microeconomics}
\author{}
\date{}
\begin{document}
\maketitle


\subsection{Definition}\label{definition}

The \textbf{Constant Elasticity of Substitution (CES)} production
function is given by:

\[F(K, L) = A\left[\alpha K^{\rho} + (1-\alpha) L^{\rho}\right]^{r/\rho}\]

where:

\begin{itemize}
\tightlist
\item
  \(A > 0\) is total factor productivity
\item
  \(\alpha \in (0,1)\) is the distribution (share) parameter
\item
  \(\rho \leq 1, \; \rho \neq 0\) governs the elasticity of substitution
\item
  \(r > 0\) governs returns to scale
\end{itemize}

The elasticity of substitution (defined and derived below) is
\(\sigma = \frac{1}{1-\rho}\). Note then
\(\rho = \frac{\sigma-1}{\sigma}\), so sometimes CES function (for
\(r=1\)) is also written as:

\[F(K, L) = A\left[\alpha K^{\frac{\sigma-1}{\sigma}} + (1-\alpha) L^{\frac{\sigma-1}{\sigma}}\right]^{\frac{\sigma}{\sigma-1}}\]

CES function nests several important special cases:

\begin{itemize}
\tightlist
\item
  Cobb-Douglas: \(\rho \to 0\) (or \(\sigma \to 1\))
\item
  Perfect substitutes: \(\rho \to 1\) (or \(\sigma \to \infty\))
\item
  Perfect complements: \(\rho \to -\infty\) (or \(\sigma \to 0\))
\end{itemize}

\subsection{Returns to Scale}\label{returns-to-scale}

Scale all inputs by \(\lambda > 0\):

\[F(\lambda K, \lambda L) = A\left[\alpha (\lambda K)^{\rho} + (1-\alpha)(\lambda L)^{\rho}\right]^{r/\rho}\]

\[= A\left[\lambda^{\rho}\left(\alpha K^{\rho} + (1-\alpha) L^{\rho}\right)\right]^{r/\rho}\]

\[= \lambda^{r} \cdot A\left[\alpha K^{\rho} + (1-\alpha) L^{\rho}\right]^{r/\rho}\]

\[\boxed{F(\lambda K, \lambda L) = \lambda^{r} \cdot F(K,L)}\]

The function is homogeneous of degree \(r\): increasing returns if
\(r > 1\), constant returns if \(r = 1\), and decreasing returns if
\(r < 1\).

\subsection{Marginal Products}\label{marginal-products}

Let \(Z = \alpha K^{\rho} + (1-\alpha) L^{\rho}\) so that
\(F = A Z^{r/\rho}\). By the chain rule:

\[MP_K = \frac{\partial F}{\partial K} = A \cdot \frac{r}{\rho} Z^{r/\rho - 1} \cdot \alpha \rho K^{\rho - 1} = A r \alpha K^{\rho-1} Z^{r/\rho - 1}\]

\[\boxed{MP_K = A r \alpha K^{\rho - 1}\left[\alpha K^{\rho} + (1-\alpha) L^{\rho}\right]^{r/\rho - 1}}\]

By symmetry \(MP_L = A r (1-\alpha) L^{\rho - 1} Z^{r/\rho - 1}\) or:

\[\boxed{MP_L = A r (1-\alpha) L^{\rho - 1}\left[\alpha K^{\rho} + (1-\alpha) L^{\rho}\right]^{r/\rho - 1}}\]

\subsection{Marginal Rate of Technical
Substitution}\label{marginal-rate-of-technical-substitution}

The MRTS is:

\[MRTS = \frac{MP_L}{MP_K} = \frac{A r (1-\alpha) L^{\rho - 1} Z^{r/\rho - 1}}{A r \alpha K^{\rho - 1} Z^{r/\rho - 1}}\]

Canceling common terms, we get:

\[\boxed{MRTS = \frac{1-\alpha}{\alpha}\left(\frac{L}{K}\right)^{\rho-1}}\]

The MRTS captures the rate at which the firm can substitute capital for
labor while keeping output constant. For example, if \(MRTS = 2\), then
the firm can can replace one unit of labor with two units of capital
without changing output.

From the expression above, we can see that the MRTS at all input ratios
is higher when:

\begin{itemize}
\tightlist
\item
  \(\alpha\) is smaller (i.e., when labor is more important in
  production, so we need more capital to substitute for labor)
\item
  \(\rho\) is smaller (i.e., when inputs are more easily substitutable)
\end{itemize}

Additionally, the capital units needed to substitute for one unit of
labor increases as the input ratio \(K/L\) increases (power \(\rho-1\)
is negative.)

\subsection{Elasticity of
Substitution}\label{elasticity-of-substitution}

The MRTS tells us the rate at which a firm can trade one input for
another while maintaining the same output level, but it doesn't tell us
how responsive the firm's input mix is to changes in that trade-off
rate. In other words, MRTS captures the slope of the isoquant at a
point, but it doesn't tell us about the shape of the isoquant and how
quickly the slope changes as we move along the isoquant. The elasticity
of substitution captures this responsiveness.

In particular, the elasticity of substitution is defined as:

\[\sigma = \frac{d\ln(K/L)}{d\ln(MRTS)}\]

Remember changes in logs capture percentage changes
(\(d\ln(x) =  dx/x\)). Therefore, \(\sigma\) captures the percentage
change in the input ratio for a given percentage change in the MRTS.

\begin{figure}

\centering{

\pandocbounded{\includegraphics[keepaspectratio]{lec3-ces-prod-func_files/figure-pdf/fig-elasticity-two-panel-output-1.pdf}}

}

\caption{\label{fig-elasticity-two-panel}Two isoquants with the same
MRTS at point A but different curvatures}

\end{figure}%

\begin{verbatim}
  Case   σ  K/L at A  MRTS at A  K/L at B  MRTS at B
High σ 2.0       2.0        2.0      0.18        0.6
 Low σ 0.5       2.0        2.0      1.09        0.6
\end{verbatim}

In the above example, both isoquants start at point A with K/L = 2.0 and
MRTS = 2.0. At point B, both have the same new MRTS of 0.6. So the
\emph{change} in the marginal rate is identical. But look at how
differently the input mix responds. With high σ, K/L plummets from 2.0
to 0.18: the firm dramatically shifts toward labor. With low σ, K/L
moves by much less, from 2.0 to 1.09.

To derive that \(\sigma = \frac{1}{1-\rho}\), we start by taking logs of
the MRTS expression:

\[\ln(MRTS) = \ln\left(\frac{\alpha}{1-\alpha}\right) + (\rho - 1)\ln\left(\frac{L}{K}\right)\]

Plugging \(\ln(L/K) = -\ln(K/L)\) and total differentiating (\(\alpha\)
and \(\rho\) are constants):

\[d\ln(MRTS) = (1 - \rho) \, d\ln(K/L)\]

Therefore:

\[\boxed{\sigma = \frac{1}{1-\rho}}\]

The elasticity of substitution is constant and it does not depend on
input levels, which gives the CES function its name.

\subsection{Cost Minimization}\label{cost-minimization}

\subsubsection{Problem}\label{problem}

For constant returns to scale (\(r = 1\)), the firm minimizes the cost
of producing output \(q\):

\[\min_{K, L} \quad wL + rK\]

subject to:

\[A\left[\alpha K^{\rho} + (1-\alpha) L^{\rho}\right]^{1/\rho} = q\]

where \(w\) is the wage and \(r\) is the rental rate of capital.

\subsubsection{Optimality Condition}\label{optimality-condition}

Setting \(MRTS = w/r\) (the price ratio):

\[\frac{1-\alpha}{\alpha}\left(\frac{L}{K}\right)^{\rho-1} = \frac{w}{r}\]

Solving for the optimal input ratio:

\[\frac{L}{K} = \left(\frac{\alpha}{1-\alpha} \cdot \frac{w}{r}\right)^{-\frac{1}{1-\rho}}\]

Or equivalently:

\[\boxed{\frac{K}{L} = \left(\frac{1-\alpha}{\alpha} \cdot \frac{r}{w}\right)^{-\frac{1}{1-\rho}} = \left(\frac{1-\alpha}{\alpha} \cdot \frac{r}{w}\right)^{-\sigma}}\]

This shows that \(\sigma\) directly governs how the input ratio responds
to changes in relative factor prices.

In applied work, the log form of the above expression is often used for
estimation:

\[\ln\left(\frac{K}{L}\right) = \underbrace{\sigma \ln\left(\frac{\alpha}{1-\alpha}\right)}_{\text{Constant}} - \sigma \ln\left(\frac{r}{w}\right)\]

So regressing log of relative input use on log of relative factor prices
yields an estimate of the elasticity of substitution \(\sigma\).

This shows that we can also express the elasticity of substitution as:

\[\sigma = \frac{d\ln(K/L)}{d\ln(r/w)}\]

where \(K/L\) is the equilibrium input ratio because in equilibrium
\(MRTS = w/r\). So another way to interpret \(\sigma\) is that it
captures the percentage change in the input ratio for a given percentage
change in relative factor prices.

\subsubsection{Conditional Factor
Demands}\label{conditional-factor-demands}

Substituting the optimal ratio back into the production constraint
yields the conditional factor demands. In particular, from the optimal
ratio we have
\(K = L \cdot \left(\frac{1-\alpha}{\alpha} \cdot \frac{r}{w}\right)^{\frac{1}{\rho-1}}\),
so we can write the production constraint as:

\[A\left[\alpha \left(L \cdot \left(\frac{1-\alpha}{\alpha} \cdot \frac{r}{w}\right)^{\frac{1}{\rho-1}}\right)^{\rho} + (1-\alpha) L^{\rho}\right]^{1/\rho} = q\]

Rearranging and solving for \(L\):
\[L(w, r, q) = \frac{q}{A}\left[\alpha \left(\frac{1-\alpha}{\alpha} \cdot \frac{r}{w}\right)^{\frac{\rho}{\rho-1}} + (1-\alpha)\right]^{-\frac{1}{\rho}}\]

Taking common denominator for the term inside the square brackets:
\[L(w, r, q) = \frac{q(\alpha w)^{\frac{1}{\rho-1}}}{A}\left[\alpha((1-\alpha)r)^{\frac{\rho}{\rho-1}} + (1-\alpha)(\alpha w)^{\frac{\rho}{\rho-1}}\right]^{-\frac{1}{\rho}}\]

Factoring out \(\alpha(1-\alpha)\) inside the square brackets (note
coming out from the brackets it will have power \(-1/\rho\))

\[L(w, r, q) = \frac{q(\alpha w)^{\frac{1}{\rho-1}}}{A(\alpha(1-\alpha))^{1/\rho}}\left[(1-\alpha)^{\frac{1}{\rho-1}}r^{\frac{\rho}{\rho-1}} + \alpha^{\frac{1}{\rho-1}} w^{\frac{\rho}{\rho-1}}\right]^{-\frac{1}{\rho}}\]

Now noting
\(K = L \cdot \left(\frac{1-\alpha}{\alpha} \cdot \frac{r}{w}\right)^{\frac{1}{\rho-1}}\),
we can write the conditional demand for capital as:

\[K(w, r, q) = \frac{q((1-\alpha)r)^{\frac{1}{\rho-1}}}{A(\alpha(1-\alpha))^{1/\rho}}\left[(1-\alpha)^{\frac{1}{\rho-1}}r^{\frac{\rho}{\rho-1}} + \alpha^{\frac{1}{\rho-1}} w^{\frac{\rho}{\rho-1}}\right]^{-\frac{1}{\rho}}\]

\subsubsection{Cost Function}\label{cost-function}

The cost function is the minimum cost of producing output \(q\) given
factor prices:

\[C(w, r, q) = wL(w, r, q) + rK(w, r, q)\]

Plugging in the expressions for \(L(w, r, q)\) and \(K(w, r, q)\):

\[C(w, r, q) = \frac{q}{A(\alpha(1-\alpha))^{1/\rho}}\left[(1-\alpha)^{\frac{1}{\rho-1}}r^{\frac{\rho}{\rho-1}} + \alpha^{\frac{1}{\rho-1}} w^{\frac{\rho}{\rho-1}}\right]^{-\frac{1}{\rho}} X\]

where
\[ X = w(\alpha w)^{\frac{1}{\rho-1}} + r((1-\alpha)r)^{\frac{1}{\rho-1}}\]

Which is the same as the term inside the square brackets, so we can
simplify to:

\[\boxed{C(w, r, q) = \frac{q}{A(\alpha(1-\alpha))^{1/\rho}}\left[(1-\alpha)^{\frac{1}{\rho-1}}r^{\frac{\rho}{\rho-1}} + \alpha^{\frac{1}{\rho-1}} w^{\frac{\rho}{\rho-1}}\right]^{\frac{\rho-1}{\rho}}}\]

Interestingly, the cost function also has a CES form in factor prices,
but with elasticity of substitution \(\rho\).




\end{document}
